\documentclass[protokoll]{dvd}

\KOMAoptions{
    headwidth = 18cm,
    footwidth = 18cm,
}

% Styrelsen
\def\ordf{Daniell Cole}
\def\vice{\vak}
\def\samo{Angelika Hagberg}
\def\kassa{Owais Tabbaa}
\def\sekr{\vak}
\def\vak{\emph{vakant}}

% mötesinfo

\def\nr{5}
\def\yr{2026}
\def\ymd{\yr-11-02}
\def\starttime{12:00}


\begin{document}

\title{Styrelsemöte \nr}
\subtitle{\yr}
\author{Styrelsen}
\date{\ymd}


\textbf{Datum:} \csname @date\endcsname\\
\textbf{Tid:} \starttime\\
\textbf{Plats:} Styrelserummet\\
\textbf{Styrelsemedlemmar:}
\begin{närvarande_förtroendevalda}
    \förtroendevald{Ordförande}{\ordf}{Ja}
    \förtroendevald{Vice Ordförande}{\vice}{---}
    \förtroendevald{Kassör}{\kassa}{Ja}
    \förtroendevald{SAMO}{\samo}{Ja}
    \förtroendevald{Sekreterare}{\sekr}{---}
\end{närvarande_förtroendevalda}

%\textbf{Övriga medlemmar:}
%\begin{närvarande_medlemmar}
%\medlem{Ingen}[]
%\end{närvarande_medlemmar}
\section{Öppnande av möte}

Mötet planeras  öppnas av \ordf\ kl \starttime

mötet öppnas 12:04

\section{Runda bordet}

Runda bordet innebär att varje person berättar hur de känner sig.
Man kan till exempel berätta att man är stressad på grund av en inlämning, irriterad på sin granne, eller bara väldigt glad därför att man ligger i fas med plugget.

\section{Formalia}

\subsection{Styrelsens beslutbarhet}

\blockquote[7 kap. 5 \S~första stycket i stadgan][]{%
    Styrelsen är endast beslutsmässig då samtliga styrelsemedlemmar har fått kallelsen till styrelsemötet och minst hälften av styrelsemedlemmarna är närvarande. Ordförande eller vice ordförande måste vara närvarande när beslut tas.
}

\blockquote[7 kap. 11 \S~första stycket i stadgan]{
Styrelsens sammanträden måste kallas minst två läsdagar innan sammanträdet. }

\blockquote[7 kap. 11 \S~sista stycket i stadgan]{
Kallelse till sammanträde får endast ske när protokollet från tidigare sammanträde är justerat och tillgängligt för medlemmarna.}

\subsubsection*{Förslag}
Kallelse skedde tekninskt sett förra mötet och förra mötets protokoll gjordes tilgängligt för styrelsen 23:e/10. Protokollet blev offentligt tillgängligt 30/10 men jag glömde annonsera. 
\begin{attsatser}
    \item Styrelsen har uppnått kraven i 7 kap. 5 § första stycket i stadgan och är därmed beslutbar.
\end{attsatser}
\subsubsection*{Beslut}
\begin{attsatser}
    \item förslaget till beslut bifalles
\end{attsatser}


\subsection{Fastställande av mötesschema}

För att styrelsen ska kunna fatta ett styrelsebeslut eller protokollföra en diskussion behöver punkten i mötesschemat där styrelsen ska fatta beslut vara inlagd eller föras in i mötesschemat senast vid den här punkten.

\subsubsection*{Förslag}

\begin{attsatser}
    \item mötesschemat fastställs utan några förändringar.
    \item lägga till diskussionspunkt på schemat gällande en incident. 
    %punkt läggs till på mötesschema
\end{attsatser}
\subsubsection*{Beslut}
\begin{attsatser}
    \item attsats 2 bifalles
\end{attsatser}

\subsection{Val av mötessekreterare}
Vi har ingen sekreterare i styrelsen men vi behöver någon som antecknar under mötets gång. 
\subsubsection*{Förslag}
\begin{attsatser}
    \item Angelika Haglund väljs till mötessekreterare. 
\end{attsatser}
\subsubsection*{Beslut}
\begin{attsatser}
    \item attsats 1 bifalles
\end{attsatser}

\subsection{Val av protokolljusterare}

Protokolljusterare har till uppgift att kontrollera att protokollet i slutändan reflekterar de faktiska besluten och diskussionerna som fördes under mötet.
Utöver protokolljusteraren så ska mötesordförande och mötessekreteraren signera protokollet.
Vid styrelsemöten ska det endast vara en justerare.
Mötesordförande och mötessekreteraren kan inte vara justerare.


\subsubsection*{Förslag}
\begin{attsatser}
    \item Owais Tabbaa väljs till protokolljusterare, 
\end{attsatser}
\subsubsection*{Beslut}
\begin{attsatser}
    \item Attsats 1 bifalles
\end{attsatser}

\section{Rapport} 
\subsection{Ordförande}
Jag har varit ganska uptagen med med tenta plugg och inlämningar men har skickat ut info om SUS-utbildning. Jag har inte fått något svar än. Sedan har jag också anmält DVD till att vara med på Kårkällar'ns rullande schema och annonserar när vi/om vi får någon onsdag. 

Jag har också fått en liten detalj-feedback på förslag 3 att förtydliga att det gäller våra interna kommittéer och inte götas kommittéer (e.g. pubbkommittéen). 


\subsection{Kassör}
Ska betala ut de senaste äskningarna någon gång i veckan. Ska fråga tidigare
kassör angående vad som ska nämnas i årsrapporten.


\subsection{SAMO}
Inget att rapportera


\newpage


\section{Diskussionspunkter}
\subsection*{Pubrunda}
Vart ska vi gå, vilken ordning? Det är en onsdag och dagen efter är vanlig läsdag, så behöver vara något vaksamma med tiden.\\ 

Börja med Kårkällar'n \textrightarrow\ Vasa brygg \textrightarrow\ Gustav vasa bar, och senare vidare mot andra lång så får vi se hur uppslutningen ser ut.

\subsection*{Incident}
Fortsättning av det som hände under Mega6 sittning under insparken. Det är
fortfarande något som tas upp och diskuteras. Det finns fortfarande frustration kvar och det
behövs pratas ut. Angelika och Daniell ska kalla in till ett möte.

\section{Besultspunkter}

\subsection{Decemberstämma}
Vi behöver ha en stämma i december, efter att alla kommittéer ha beslutat om sina nya ordförande. Kanske v. 50, andra veckan i december, 07-11/12?
\subsubsection{Beslut}
\begin{attsatser}
    \item Vi siktar på torsdag 10:e/12 vecka 50. Daniell bokar lokal och skickar ut kallelse när salsbokning är bekräftad. 
\end{attsatser}


\subsection{Stadgaändringsförslag }

\subsubsection*{Förslag 1}
Minimum 3 medlemmar i styrelsen: 
\begin{attsatser}
    \item i början av 7 kap. 3 § i stadgan införa 
    \begin{quote}
        Styrelsen kan som minst bestå av tre medlemmar varav en skall vara divisionsordförande och en divisionskassör. 
    \end{quote}
\end{attsatser}

\subsubsection*{Förslag 2}
Krav från Göta, styrelsemedlem behöver vara ordinarie medlem i Göta studentkår:
\begin{attsatser}
    \item ändra första stycket i 7 kap. 2 § från
    \begin{quote}
        Styrelsens medlemmar väljs av divisionsstämman. Mandatperioden är densamma som räkenskapsåret. Valbar till styrelsen är medlem i divisionen.
    \end{quote}
    till
    \begin{quote}
        Styrelsens medlemmar väljs av divisionsstämman. Mandatperioden är densamma som räkenskapsåret. Valbar till styrelsen är medlem i divisionen med ordinarie medlemskap i Göta studentkår.
    \end{quote}
\end{attsatser}

\subsubsection*{Förslag 3}
Ett till krav från Göta, sammanfatta vilka rättigheter föreningsmedlemmar har:
\begin{attsatser}
    \item införa i 4 kap, 4 § av stadgan 
    \begin{quote}
        Medlemskap medför: 
        \begin{itemize}
            \item rätten att kalla till divisionsstämma (6 kap, 1 §),
            \item närvaro-, yttrande, förslags- och rösträtt på divisionsstämman (6 kap, 5 §),
            \item rätt att bli medlem i styrelsen (7 kap, 2 §),
            \item närvarorätt på styrelsemöten (7 kap, 6 §),
            \item samt rätten att bli medlem i divisionens interna kommittéer (8 kap, 4 §).
        \end{itemize}
    \end{quote}
\end{attsatser}

\subsubsection*{Förslag 4}
Den stora jag främst vill ha feedback på, ny roll i styrelsen; Styrelseledamot:
\begin{attsatser}
    \item ändra 7 kap, 1§ från:
    \begin{quote}
    Styrelsen består av
		\begin{itemize}
			\item ordförande (även kallad divisionsordförande);
			\item vice ordförande;
			\item sekreterare;
			\item kassör (även kallad divisionsskassör);
			\item SAMO; och
            \end{itemize}
    \end{quote}
    till:
    \begin{quote}
        Styrelsens ordinarie medlemmar består av 
		\begin{itemize}
			\item ordförande (även kallad divisionsordförande);
			\item vice ordförande;
			\item sekreterare;
			\item kassör (även kallad divisionsskassör);
			\item SAMO; och
            \end{itemize}
            ytterligare styrelsemedlemmar består av
            \begin{itemize}
                \item ledamöter (även kallat    styrelseledamöter).
            \end{itemize}
    \end{quote}
    
    \item införa i 7 kap. 2 §:
    \begin{quote}
        Ordinarie styrelsemedlem kan inte samtidigt vara styrelseledamot.
    \end{quote}
    
    \item uppdatera 7 kap. 3 § från:
    \begin{quote}
        Styrelsen kan högst bestå av fem medlemmar.
    \end{quote}
    till: 
    \begin{quote}
        Styrelsen kan högst bestå av fem ordinarie medlemmar. 
    \end{quote}
    
    \item 
    ändra första stycket i 7 kap. 5 § från:
    \begin{quote}
        Styrelsen är endast beslutsmässig då samtliga styrelsemedlemmar har fått kallelsen till styrelsemötet och minst hälften av styrelsemedlemmarna är närvarande.
        Ordförande eller vice ordförande måste vara närvarande när beslut tas.
    \end{quote}
    till:
    \begin{quote}
        Styrelsen är endast beslutsmässig då samtliga styrelsemedlemmar har fått kallelsen till styrelsemötet och minst hälften av styrelsens ordinarie medlemmar är närvarande.
        Ordförande eller vice ordförande måste vara närvarande när beslut tas.
    \end{quote}
\end{attsatser}

Jag har inte hunnit visa förslagen till Albin men vill fortfarande ha hans åsikt om de innan vi lägger till det som förslag till stämman. 

- \ordf
\subsubsection{Förslag till beslut}
\begin{attsatser}
    \item Punkten borsdläggs tills nästa möte. 
\end{attsatser}
\subsubsection{Beslut} 
\begin{attsatser}     
\item %Attsats 1 bifalles. 
\end{attsatser}


\subsection{Äskning: ConCats; Spelkväll}
\begin{quote}
Hej!
    
Vi i ConCats tänker hålla ett spelkväll den 5:e november, vi förväntar runt 10-15 personer så vi vill äska 400 kr för mat till dem som kommer till eventet.
    
Mvh Owais "Mirage" Tabbaa
\end{quote}

\subsubsection*{Förslag till beslut}
\begin{attsatser}
    \item Att bevilja äskan för 400 kr till fika 05/11. 
\end{attsatser}
\subsection{Beslut}
\begin{attsatser}
    \item attsats 1 bifalles
\end{attsatser}

\subsection{Äskning: }
\begin{quote}
Hej!
DVRK vill ha 300kr till fika till vårt aspevent som vi ska hålla den 6/11. Vi ska hålla i en frågesport och vi estimerar att det är runt 13 personer som kommer inkl DVRK medlemmarna.
Vänligen,
Halimo "Lost\&\&Found" Hassan. 
\end{quote}

\subsubsection*{Förslag till beslut}
\begin{attsatser}
    \item Att bevilja äskan för 300 kr till fika 06/11. 
\end{attsatser}
\subsection{Beslut}
\begin{attsatser}
    \item attsats 1 bifalles
\end{attsatser}



\section{Per Capsulam-beslut}
\blockquote[]{
När per capsulam-beslut har fattats måste beslutet dokumenteras,\\
\ \cdot\ \ på divisionens Discordserver som är tillgänglig alla studenter;\\
\ \cdot\ \ i protokollet för det första styrelsesammanträdet efter beslut har fattats.\\

Dokumentationen ska innehålla\\
\ \cdot\ \ vilka att-satser som har bifallits;\\
\ \cdot\ \ omröstningsperioden.\\

Om dessa krav inte uppfylls upphävs beslutet.
}

\textit{Inga per capsulam-beslut har tagits}

\newpage

\section{Avslutande av möte}

\subsection{Nästa möte}
2026-11-17 - 12:00

\subsection{Mötets avslutande}
Mötet avslutades kl 13:01 

\styrelsesignaturer

\end{document}